% EDITING: type inside the final {...} of each field. Change the shown font size
% (for example, font=10pt) in that same field block when needed.

\GrantSection{2}{Project Summary, Public Narrative, and Nontechnical Overview}

\GrantGuidance{For the NIH baseline, keep the Project Summary/Abstract to approximately 30 lines and the Project Narrative to three sentences. Other funders may impose word or character limits instead. The character caps below are conservative drafting controls and can be changed in the section source.}

\GrantLineField[id=summary-project-title,font=11.5pt,maxchars=220]{Project title}{Daraxonrasib Phase 1 LLM-Directed Robotic Whipple in KRAS-Mutated PDAC}
\GrantLineField[id=summary-condition,font=12pt,maxchars=180]{Primary disease, condition, population, or scientific area}{KRAS-Mutated Pancreatic Ductal Adenocarcinoma}
\GrantLineField[id=summary-science-areas,font=12pt,maxchars=80]{Primary and secondary scientific-area codes, if applicable}{3 CTR, 5 IE}

\GrantTextArea[
  id=project-abstract, font=9.5pt, height=4.3in, maxchars=3200
]{Project Summary / Abstract (approximately 30 lines)}{This project advances a Phase 1, first-in-human, prospective, open-label, single-arm clinical study of perioperative daraxonrasib (RMC-6236) administered with an on-premises LLM-directed robotic pancreaticoduodenectomy for adults with resectable or borderline-resectable KRAS G12-mutated pancreatic ductal adenocarcinoma. The major scientific premise is that a targeted perioperative RAS(ON) inhibitor paired with a tightly governed robotic procedure and a verification-first computational control layer can be translated into a clinically auditable early-phase program without relaxing established standards for patient protection, investigator oversight, or data integrity \citep{phase1ind,onpremwhippl,cfr312paiuni,cfr050paiuni,ichpaistduni}. The project addresses a central translational gap in oncology: drug development, robotic systems, human-subject protections, and trial operations are usually constructed as separate handoffs, ^^J whereas this program treats them as one prospectively specified investigational system. The public foundation for the application includes the IND package, Phase 1 and Phase 2 protocol records, repository-based LLM methods for synchronized preparation of regulatory and operational documents, and PDAC simulation studies spanning large-scale in silico trials, quantitative systems pharmacology, and digital-twin frameworks that inform design assumptions, scenario testing, and verification planning without substituting for human clinical judgment \citep{phase1ind,phase1guidance,onpremwhippl,daraxwhipple,oncotrialllm,oncotrialeff,qspmetpancre,chatgpt100kp,pdacdigtwinp,fdadigtwinpc}. The immediate clinical purpose is not ^^J to prove comparative efficacy, but to estimate safety, tolerability, perioperative feasibility, and operational readiness strong enough to justify later comparative testing. Primary outputs will include dose and procedure feasibility estimates, structured safety data, robotic and cyber-physical performance logs, implementation lessons for regulated LLM-supported trial conduct, and a reusable regulatory and data-governance template for future Physical AI oncology investigations. If successful, the work would define a practical path for responsibly evaluating integrated drug-device-AI interventions in severe malignancy; if unsuccessful, it will still generate valuable stopping criteria, risk controls, and methodological lessons for the field \citep{clintrustpai,vvuqllmverif,h2pancrevvuq,natlpaiplatf,paiautospons}.
}
 

 
\GrantTextArea[
  id=project-narrative, font=10pt, height=1.15in, maxchars=900
]{Project Narrative / Public Health Relevance (three plain-language sentences)}{Pancreatic cancer remains difficult to treat, and early-phase trials require extensive scientific, regulatory, and operational preparation before patients can benefit. This project combines a targeted KRAS-directed drug strategy with an LLM-directed robotic surgical concept and uses computational trial-development workflows to improve the efficiency and auditability of trial preparation. Funding would support ^^Jresponsible early-phase clinical translation of this integrated pancreatic cancer trial program. \citep{phase1ind,phase1guidance,onpremwhippl,qspmetpancre}
}
\GrantTextArea[
  id=project-lay-summary, font=10pt, height=3.23in, maxchars=1900
]{Nontechnical overview for a foundation, patient organization, donor, or corporate review committee}{Pancreatic ductal adenocarcinoma is a highly lethal cancer, and patients who are candidates for surgery still face major risks from disease recurrence, difficult operations, and delays in translating new therapies into carefully controlled clinical studies. This project proposes an early-phase study that combines perioperative daraxonrasib, a RAS(ON) multi-selective inhibitor, with an on-premises large language model-directed robotic Whipple procedure for patients with resectable or borderline-resectable KRAS G12-mutated disease \citep{phase1ind,onpremwhippl}. ^^J ^^JThe project does not assume that artificial intelligence or robotics are safe merely because they can generate an action. Every patient-facing action would be subject to predefined human authorization, ^^Jaudit logging, emergency-stop controls, and verification, validation, and uncertainty-quantification gates that accept, escalate, or block the action before execution \citep{clintrustpai,vvuqllmverif,h2pancrevvuq}. ^^J ^^JFunding would support the regulatory, clinical, engineering, data, and safety work needed to determine whether this combined drug-device approach can be studied responsibly in humans. The immediate ^^Jgoal is not to claim efficacy, but to establish safety, feasibility, dose and procedure readiness, and a transparent evidence base for a later comparative trial \citep{phase1ind,daraxwhipple}
}

\GrantSubsection{One-line decision summary}
\GrantTextArea[id=project-decision-summary,font=10pt, height=0.6in,maxchars=500]{In one sentence: what will the requested funding make possible that cannot otherwise happen?}{Funding will enable a governed first-in-human evaluation of perioperative daraxonrasib with an LLM-directed robotic Whipple, including the regulatory, clinical, safety, and VVUQ infrastructure that simulation alone cannot establish \citep{phase1ind,onpremwhippl,vvuqllmverif}.}

